\subsection{药物辅助性侵（"迷奸"）：识别、应急与法律}

前文在约会软件风险与同意章节中多次提到"被下药"，但都只是一句话。本节把它作为独立的实操主题讲完——不是因为这件事常见到人人会遇到，而是因为它的应对有严格的时间窗，事前知道和事后才知道，后果完全不同。

\textbf{一、常见药物与识别}
\begin{itemize}
\item \textbf{GHB（γ-羟丁酸，俗称"G 水"）}：无色无味、起效 10-20 分钟、持续数小时。典型表现是\textbf{突发嗜睡、意识模糊、肌肉张力骤降}，与饮酒量不成比例；剂量与效果个体差异极大，稍多即致昏迷与呼吸抑制；
\item \textbf{苯二氮䓬类}（氟硝西泮/Rohypnol"罗眠乐"、三唑仑等）：镇静、顺行性遗忘明显——受害者常表现为"断片"，对事件无法回忆，这也正是取证与陈述困难的来源；
\item \textbf{酒精本身}：最常见也最被低估。法律意义上的关键在于\textbf{失能（incapacitation）}——当一个人因酒精或药物失去理解与表达同意的能力时，其"同意"无效。醉酒者的沉默、配合或半昏迷状态都不是同意；
\item \textbf{可识别的信号}：饮品出现异常咸味或苦味（GHB 有咸味）、喝得不多却突然极度困倦、与既往饮酒表现明显不符、记忆出现空白段。出现这些情况，第一反应应该是\textbf{求助与离开}。
\end{itemize}

\textbf{二、当事人应急（时间敏感）}
\begin{itemize}
\item \textbf{优先保安全}：联系可信任的人来接，不要独自离开，不要接受"我送你回家"；
\item \textbf{保留证据}：\textbf{不洗澡、不刷牙、不更换衣物、不冲洗下身}；尽量在 72 小时内（越早越好）到医院急诊或专门的性侵取证机构进行\textbf{药物检测与法医取证}——GHB 在血液中 8-12 小时即难以检出、尿液中约 12 小时，过了窗口就无法补证；
\item \textbf{医学处理}：72 小时内评估 HIV 阻断（PEP）、紧急避孕（含米非司酮方案）、性传播感染预防性治疗与乙肝疫苗接种。这一步与"是否报警"无关，先做医学处置；
\item \textbf{心理支持}：急性应激反应是正常的，不是"软弱"。国内有性侵受害者援助热线与心理危机干预热线，尽早联系可显著降低后续创伤后应激障碍的风险。
\end{itemize}

\textbf{三、法律要点}
\begin{itemize}
\item 使用药物使他人失去反抗能力后发生性行为，构成\textbf{强奸}——是使用"其他手段"实施的强奸，量刑上属从重情节；
\item \textbf{同意能力}是判定核心：严重醉酒、昏迷、药物影响下的"同意"在法律上无效；未明确同意的性行为即构成性侵，无需存在暴力反抗痕迹；
\item \textbf{举证现实}：此类案件取证困难，及时报案与保留证据（就医记录、检测结果、聊天记录、场所监控）是维权的关键。怀疑被下药时，请务必就医而不是自行"过去就算了"——身体与法律两个层面都留着可用的路径。
\end{itemize}
